\section{Secuencias largas de proteínas: Performer y Protein LM}

En la sección clínica, los retos principales eran el conocimiento, la evaluación y el riesgo; al pasar a las proteínas, el primer obstáculo se vuelve la longitud de la secuencia y la complejidad computacional. Una secuencia de aminoácidos puede darse como entrada al modelo igual que una secuencia de tokens, pero las proteínas largas hacen que la matriz de correlaciones $n\times n$ de la attention estándar consuma rápidamente el cómputo y la memoria. Aquí, Performer no es un recurso abstracto de eficiencia, sino la condición previa para poder entrenar con secuencias biológicas más largas.

\subsection{Por qué las proteínas necesitan contexto largo}

La diapositiva de transición de sección cambia la perspectiva del lenguaje de los médicos al lenguaje de las moléculas. La siguiente diapositiva señala que la longitud de las proteínas varía en un intervalo muy amplio, y que aminoácidos lejanos en la secuencia pueden quedar próximos entre sí después del plegamiento y determinar juntos el sitio activo. Si solo se observa una ventana pequeña y fija, se puede perder esta dependencia; si se usa directamente la self-attention global, hay que pagar una complejidad cuadrática.

\lecturefigure{slide-44-section-proteins.jpg}{Cambio de sección: aplicaciones en proteínas}{Diapositiva V044 recuperada de la grabación oficial; Stanford CS25 Lecture 17}

\lecturefigure{slide-45-long-biological-sequences.jpg}{Problemas de modelado y de cómputo que plantean las secuencias biológicas largas}{Diapositiva V045 recuperada de la grabación oficial; Stanford CS25 Lecture 17}

La scaled dot-product attention estándar es
$$
\operatorname{Attn}(Q,K,V)=
\operatorname{softmax}\!\left(\frac{QK^{\top}}{\sqrt d}\right)V.
$$
Donde $Q,K,V\in\mathbb{R}^{n\times d}$ son, respectivamente, query, key y value; $n$ es la longitud de la secuencia y $d$ es la dimensión de cada head. El cuello de botella proviene de $QK^{\top}\in\mathbb{R}^{n\times n}$: aunque al final solo se necesite una salida de $n\times d$, hay que procesar, de forma explícita o implícita, todos los pares de posiciones.

\begin{warningbox}{La dificultad de las secuencias largas no es solo «falta de memoria de GPU»}
La attention cuadrática aumenta a la vez el tiempo de entrenamiento, las activaciones intermedias y la presión de comunicación. El truncamiento simple ahorra cómputo, pero puede eliminar motifs lejanos o relaciones entre dominios estructurales; la attention por bloques, a su vez, puede separar posiciones que en realidad están relacionadas. El objetivo de los métodos de eficiencia es lograr un compromiso controlado entre el presupuesto de cómputo y la información global.
\end{warningbox}

\subsection{Performer: reordenar el cálculo con kernel features}

La subsección anterior ya ubicó el cuello de botella en la score matrix completa de $n\times n$; esta sección responde, además, cómo cambiar el orden del cálculo sin fragmentar la secuencia. La diapositiva de Performer presenta la idea de FAVOR+: aproximar el kernel softmax como un producto interno de características, de modo que las dependencias globales todavía puedan propagarse a través de un estadístico comprimido. Al leer la figura hay que observar al mismo tiempo el reordenamiento matemático, el error de la aproximación aleatoria y las curvas en hardware real, en lugar de tomar «linear attention» solo como un lema de complejidad.
$$
\exp(q^{\top}k)\approx \phi(q)^{\top}\phi(k),
$$
y después, por asociatividad, se agrega primero $\phi(K)^{\top}V$, en lugar de construir la attention matrix completa. Si se omiten los detalles de normalización, la salida de una sola query puede escribirse como
$$
\operatorname{Attn}(q_i,K,V)\approx
\frac{\phi(q_i)^{\top}\left(\sum_j \phi(k_j)v_j^{\top}\right)}
{\phi(q_i)^{\top}\left(\sum_j \phi(k_j)\right)}.
$$
Donde $\phi(\cdot)$ es un mapeo aleatorio de características positivas, y $q_i,k_j,v_j$ representan, respectivamente, la $i$-ésima query y el $j$-ésimo key/value. Después del reordenamiento ya no se guardan $n^2$ pairwise scores, y el tiempo y el espacio pueden crecer de forma aproximadamente lineal con $n$; el costo es la varianza que introduce la kernel approximation y la complejidad de implementación.

\lecturefigure{slide-46-performer-kernel-attention.jpg}{Performer: aproximación de la softmax attention con características aleatorias}{Diapositiva V046 recuperada de la grabación oficial; Stanford CS25 Lecture 17}

\lecturefigure{slide-47-performer-speed-complexity.jpg}{Comparación de velocidad, memoria y complejidad de Performer}{Diapositiva V047 recuperada de la grabación oficial; Stanford CS25 Lecture 17}

\begin{lstlisting}[caption={Orden de cálculo de la attention linealizada (esquema)}]
def performer_attention(query, key, value, phi):
    q_features = phi(query)
    k_features = phi(key)
    kv_summary = k_features.T @ value
    k_summary = k_features.sum(axis=0)
    numerator = q_features @ kv_summary
    denominator = q_features @ k_summary
    return numerator / denominator[:, None]
\end{lstlisting}

\begin{knowledgebox}{Lectura de la figura: cómo leer la gráfica de complejidad}
Primero se fijan la dimensión del modelo y el lote, y luego se aumenta la sequence length a lo largo del eje horizontal; las curvas de tiempo y de memoria de la attention estándar suben más rápido, mientras que el crecimiento de Performer se acerca más a lo lineal. La figura respalda que «las secuencias largas son más escalables», pero no demuestra que sea más rápido para todas las longitudes, todo el hardware y todos los objetivos de precisión; el número de características aleatorias, la implementación del kernel y los términos constantes modifican el punto de cruce.
\end{knowledgebox}

Desde el punto de vista de la implementación, la ganancia de la attention lineal depende de que los tensores intermedios se calculen realmente en el orden reordenado. Si, para depurar, el código todavía genera explícitamente los pairwise scores, la complejidad teórica no se materializa; si el número de features aleatorias es demasiado grande, el término constante de $n r d$ también puede superar a una attention estándar bien optimizada. En cambio, en una tarea como la de proteínas, con una distribución de longitudes muy amplia, puede elegirse la implementación por bucket de longitud: usar attention exacta para secuencias cortas, cambiar a Performer para las largas y comparar la salida aproximada con la exacta en un conjunto de muestras que el presupuesto permita. Esta validación vincula «más rápido» con «sesgo aceptable» y evita medir solo el rendimiento ignorando los cambios en el comportamiento del modelo.

\teachervoice{El ponente coloca Performer al inicio de la sección de proteínas porque las secuencias biológicas chocan con el costo de la attention cuadrática «antes de empezar a hacer razonamiento biológico». Lo importante en la clase no es memorizar el nombre FAVOR+, sino ver que la distribución de longitudes de los datos del dominio, a su vez, impulsa nuevos attention algorithms.}

\subsection{Qué aprende un Protein LM}

La diapositiva de resultados muestra la ganancia de los modelos de lenguaje en protein sequence modeling, y la diapositiva de attention visualization dibuja la atención entre distintos amino acids. Una interpretación razonable es la siguiente: el preentrenamiento obliga al modelo a aprovechar motifs conservados, el entorno químico local y las relaciones entre residuos más lejanos para predecir las posiciones enmascaradas, por lo que en las representaciones aparecen patrones relacionados con la estructura o la función.

\lecturefigure{slide-48-protein-lm-results.jpg}{Resultados experimentales de protein language modeling}{Diapositiva V048 recuperada de la grabación oficial; Stanford CS25 Lecture 17}

\lecturefigure{slide-49-amino-acid-attention.jpg}{Attention pattern entre aminoácidos y dependencias complejas}{Diapositiva V049 recuperada de la grabación oficial; Stanford CS25 Lecture 17}

\begin{warningbox}{Un attention map no es una explicación causal}
Un attention weight alto indica que ciertas posiciones interactúan con fuerza en el cálculo del modelo; no equivale a que ese par de residuos esté en contacto directo en la proteína real, ni demuestra que provoquen alguna función. Una conclusión biológica confiable requiere además datos estructurales, mutagenesis u otros experimentos. La visualización sirve sobre todo para plantear hipótesis y revisar el modelo, no para sustituir la validación experimental.
\end{warningbox}

\begin{importantbox}{El sentido del caso Performer en el curso}
Cuando los datos son largos por naturaleza, la arquitectura del modelo y la complejidad del algoritmo determinan qué preguntas científicas pueden plantearse. Performer amplía el contexto procesable mediante una kernel attention aproximada; el protein LM, por su parte, aprovecha ese contexto para aprender regularidades estadísticas. El avance en eficiencia y el descubrimiento en el dominio son dos niveles de evidencia y no deben fusionarse en una sola conclusión.
\end{importantbox}

\subsection{Resumen de la sección}

Las secuencias de proteínas exponen de manera muy directa el costo de contexto largo del Transformer. Performer evita la attention matrix completa de $n\times n$ mediante características aleatorias y la asociatividad, lo que hace más viable el entrenamiento con secuencias largas; las representaciones y los attention patterns del protein LM pueden aportar indicios biológicos, pero todavía requieren validación estructural y experimental.
